\chapter{The Standard Model as a Hybrid Quantum Field Theory}
\label{ch:standard-model-hybrid}
\ConventionsLorentzian

\section{Gauge and Matter Data}

\begin{convention}[Electroweak charge and trace normalizations]
\label{conv:sm-electroweak-trace-normalization}
The hypercharge coordinate is normalized by
\[
  Q_{\rm em}=T^3_L+Y ,
\]
where \(T^3_L\) is the third generator of the weak \(SU(2)_L\) doublet in the
usual electroweak normalization.  For transparent mass and charge formulae,
this chapter temporarily uses the conventional explicit-coupling fields
\(\mathcal G_\mu,\mathcal W_\mu,\mathcal B_\mu\).  The conversion to the
monograph's Yang--Mills convention of
Chapter~\ref{ch:vol2-classical-yang-mills-matter} is
\[
  A^{(3)}_\mu=g_3\mathcal G^A_\mu T^A,\qquad
  A^{(2)}_\mu=g_2\mathcal W^a_\mu \tau^a,\qquad
  A^{(1)}_\mu=g_1\mathcal B_\mu Y ,
\]
where \(T^A\) and \(\tau^a=\sigma^a/2\) are the conventional half-trace
generators.  If \(t^a=\sqrt2\,\tau^a\) denotes the trace-delta basis used in
Chapter~\ref{ch:vol2-classical-yang-mills-matter}, then
\(g_{2,\rm mono}=g_2/\sqrt2\) in a matter vertex.  The invariant kinetic
terms are equivalently written as explicit-coupling kinetic terms for
\(\mathcal G,\mathcal W,\mathcal B\) or as
\(-\frac1{4g^2}\operatorname{tr}(F_{\mu\nu}F^{\mu\nu})\) for the
coupling-absorbed connections.
\end{convention}

\begin{definition}[One-generation left-Weyl representation]
\label{def:sm-one-generation-left-weyl}
For one generation, written as left-handed Weyl fields, the representation is
\[
\begin{array}{c|c|c}
  \text{field} & SU(3)_c\times SU(2)_L & Y \\ \hline
  Q_L=(u_L,d_L) & (\mathbf 3,\mathbf 2) & 1/6\\
  u^c & (\overline{\mathbf 3},\mathbf 1) & -2/3\\
  d^c & (\overline{\mathbf 3},\mathbf 1) & 1/3\\
  L_L=(\nu_L,e_L) & (\mathbf 1,\mathbf 2) & -1/2\\
  e^c & (\mathbf 1,\mathbf 1) & 1 .
\end{array}
\]
Here \(u^c,d^c,e^c\) are left-handed charge-conjugate fields, not
right-handed fields in the same representation.  A gauge-singlet neutrino
field, if included in the same left-handed convention as \(\nu^c\), has
representation \((\mathbf1,\mathbf1)_0\) and does not affect the local gauge
anomaly sums.
\end{definition}

\begin{definition}[Minimal Higgs and Yukawa datum]
\label{def:sm-higgs-yukawa-datum}
The minimal Higgs scalar is
\[
  H\in(\mathbf1,\mathbf2)_{1/2}.
\]
Set
\[
  \widetilde H:=\ii\sigma^2H^*,
  \qquad
  \widetilde H\in(\mathbf1,\mathbf2)_{-1/2},
\]
where the second statement uses the pseudoreality of the \(SU(2)\) doublet.
In the left-handed Weyl convention of
Definition~\ref{def:sm-one-generation-left-weyl}, the renormalizable Yukawa
sector is
\[
  \mathcal L_Y
  =
  - (Y_u)_{IJ}\,(Q_L^I H)\,u^{cJ}
  - (Y_d)_{IJ}\,(Q_L^I\widetilde H)\,d^{cJ}
  - (Y_e)_{IJ}\,(L_L^I\widetilde H)\,e^{cJ}
  + \text{Hermitian conjugate},
\]
where \(I,J=1,2,3\) are generation labels and all gauge indices are contracted
by the invariant tensors of \(SU(3)_c\) and \(SU(2)_L\).  In components,
\((QH)u^c\) abbreviates the contraction of the two Weyl spinor indices by
\(\epsilon_{\alpha\beta}\), the two weak doublet indices by
\(\epsilon_{ij}\), and the color indices by the natural
\(\mathbf3\otimes\overline{\mathbf3}\to\mathbf1\) pairing.  The displayed
\(Y_u,Y_d,Y_e\) are complex \(3\times3\) matrices of local coupling
coordinates.
\end{definition}

\begin{definition}[Local Standard Model gauge datum]
\label{def:local-standard-model-gauge-datum}
The local Standard Model gauge algebra is
\[
  \mathfrak g_{\rm SM}
  =
  \mathfrak{su}(3)_c\oplus\mathfrak{su}(2)_L\oplus\mathfrak u(1)_Y .
\]
The covering compact group is
\[
  \widetilde G_{\rm SM}=SU(3)_c\times SU(2)_L\times U(1)_Y .
\]
A global form is a quotient
\[
  G_{\rm SM}=\widetilde G_{\rm SM}/\Gamma ,
  \qquad
  \Gamma\subset Z(\widetilde G_{\rm SM}),
\]
for which the local fields in
Definitions~\ref{def:sm-one-generation-left-weyl} and
\ref{def:sm-higgs-yukawa-datum} descend to representations of
\(G_{\rm SM}\).  On a spin four-manifold \(M\), a finite-cutoff local Standard
Model datum consists of a principal \(G_{\rm SM}\)-bundle \(P\to M\), a
connection on \(P\), three copies of the chiral matter representation in
Definition~\ref{def:sm-one-generation-left-weyl}, one Higgs scalar in
Definition~\ref{def:sm-higgs-yukawa-datum}, and the local action functional
obtained from the Yang--Mills, scalar, Yukawa, and gauge-fixing/BV data after
a regulator has been specified.
\end{definition}

The quotient \(\Gamma\) is invisible in the Lie algebra and in local
perturbative vertices, but it changes the allowed bundles, line operators,
and higher-form backgrounds.  Those global questions belong to the
generalized-symmetry analysis in
Chapter~\ref{ch:global-forms-higher-form-symmetry}.  This chapter uses only
the local representation-theoretic data unless a bundle or line operator is
explicitly mentioned.

\begin{proposition}[Gauge invariance of the minimal Yukawa datum]
\label{prop:sm-yukawa-gauge-invariance}
The three terms in Definition~\ref{def:sm-higgs-yukawa-datum} are precisely
the gauge-invariant local monomials that are bilinear in the listed
left-handed Weyl fields, linear in \(H\) or \(\widetilde H\), and of
engineering dimension at most four, up to flavor tensors and Hermitian
conjugation.
\end{proposition}

\begin{proof}
Lorentz invariance forces the two left-handed Weyl spinor indices in each
fermion bilinear to be contracted by the antisymmetric invariant tensor of
\(SL(2,\mathbb C)\).  Color invariance requires \(Q_L\) to pair with either
\(u^c\) or \(d^c\), because the quark singlet contraction is
\(\mathbf3\otimes\overline{\mathbf3}\to\mathbf1\); the lepton bilinear is
already color neutral.  Weak \(SU(2)\)-invariance then forces the remaining
weak doublet index to be contracted with either \(H\) or \(\widetilde H\) by
the invariant tensor of the doublet representation.

The hypercharge sums determine which Higgs doublet occurs:
\[
  \frac16+\frac12-\frac23=0,\qquad
  \frac16-\frac12+\frac13=0,\qquad
  -\frac12-\frac12+1=0 .
\]
Thus \(Q_LHu^c\), \(Q_L\widetilde H d^c\), and
\(L_L\widetilde H e^c\) are gauge invariant.  The opposite choices of
\(H\) and \(\widetilde H\) have nonzero hypercharge, hence do not define
local gauge-invariant monomials in the left-handed convention.  Power
counting excludes additional renormalizable matter monomials with more fields
or derivatives in this minimal field content.  Flavor indices are spectators
for the gauge group, so the allowed coefficients are the three arbitrary
complex flavor matrices shown above.
\end{proof}

Flavor rotations of the fermion fields remove unphysical parameters; the
remaining charged-current quark mixing is encoded in the CKM matrix.  If
neutrino masses are included through dimension-five or singlet-neutrino
operators, a separate PMNS mixing matrix appears in the lepton sector.

\section{Electroweak Breaking and Mass Coordinates}

\begin{definition}[Electroweak scalar chart]
\label{def:sm-electroweak-scalar-chart}
The minimal renormalizable electroweak scalar potential is
\[
  V(H)=\lambda\left(H^\dagger H-\frac{v^2}{2}\right)^2 ,
  \qquad \lambda>0 .
\]
In the local field-space chart used below,
\[
  H(x)=\frac1{\sqrt2}
  \begin{pmatrix}
    \chi_1(x)+\ii\chi_2(x)\\
    v+h(x)+\ii\chi_3(x)
  \end{pmatrix}.
\]
The reference vacuum vector is
\[
  H_0=\frac1{\sqrt2}\begin{pmatrix}0\\ v\end{pmatrix}
\]
in the same chart.  The fields \(\chi_1,\chi_2,\chi_3\) are local
coordinates tangent to the broken gauge orbit at \(H_0\), while \(h\) is the
radial scalar coordinate.
\end{definition}

\begin{proposition}[Electromagnetic stabilizer and weak-boson masses]
\label{prop:sm-electroweak-stabilizer-masses}
In the convention of
Convention~\ref{conv:sm-electroweak-trace-normalization}, the stabilizer of
\(H_0\) inside the local group \(SU(2)_L\times U(1)_Y\) is generated by
\[
  Q_{\rm em}=T^3_L+Y .
\]
The quadratic gauge-boson mass coordinates are
\[
  W^\pm_\mu=\frac1{\sqrt2}(\mathcal W^1_\mu\mp\ii\mathcal W^2_\mu),
  \qquad
  Z_\mu=\frac{g_2\mathcal W^3_\mu-g_1\mathcal B_\mu}
              {\sqrt{g_1^2+g_2^2}},
  \qquad
  A_\mu=\frac{g_1\mathcal W^3_\mu+g_2\mathcal B_\mu}
              {\sqrt{g_1^2+g_2^2}},
\]
with
\[
  m_W=\frac{g_2v}{2},
  \qquad
  m_Z=\frac{v}{2}\sqrt{g_1^2+g_2^2},
  \qquad
  m_A=0 .
\]
The coupling of the massless field is
\[
  e=g_2\sin\theta_W=g_1\cos\theta_W,
  \qquad
  \tan\theta_W=\frac{g_1}{g_2}.
\]
\end{proposition}

\begin{proof}
The weak generator \(T^3_L=\operatorname{diag}(1/2,-1/2)\) acts on
\(H_0\) by \(-H_0/2\), while the hypercharge generator acts by
\(H_0/2\).  Hence \((T^3_L+Y)H_0=0\).  The generators \(T^1_L\) and
\(T^2_L\) move \(H_0\) into the upper component, and a general combination
\(aT^3_L+bY\) annihilates \(H_0\) only when \(a=b\).  Therefore the
one-dimensional stabilizer Lie algebra is spanned by \(Q_{\rm em}\).

The covariant derivative in the explicit-coupling convention is
\[
  D_\mu H
  =
  \left(
    \partial_\mu-\ii g_2\mathcal W^a_\mu\tau^a-\ii g_1\mathcal B_\mu Y
  \right)H .
\]
Using
\[
  \tau^1H_0=\frac{v}{2\sqrt2}\begin{pmatrix}1\\0\end{pmatrix},
  \qquad
  \tau^2H_0=\frac{v}{2\sqrt2}\begin{pmatrix}-\ii\\0\end{pmatrix},
  \qquad
  \tau^3H_0=-\frac12H_0,
  \qquad
  YH_0=\frac12H_0 ,
\]
the quadratic term from \((D_\mu H)^\dagger(D^\mu H)\) at \(H=H_0\) is
\[
  (D_\mu H_0)^\dagger(D^\mu H_0)
  =
  \frac{v^2}{8}
  \left[
    g_2^2(\mathcal W^1_\mu\mathcal W^{1\mu}
          +\mathcal W^2_\mu\mathcal W^{2\mu})
    +(-g_2\mathcal W^3_\mu+g_1\mathcal B_\mu)^2
  \right].
\]
The charged part is
\((g_2^2v^2/4)W^+_\mu W^{-\mu}\), so
\(m_W^2=g_2^2v^2/4\).  The neutral quadratic form has rank one: the massive
unit vector is \(Z_\mu\), the null vector is \(A_\mu\), and comparison with
the real-vector mass term \(\frac12m_Z^2Z_\mu Z^\mu\) gives
\(m_Z^2=(g_1^2+g_2^2)v^2/4\).

Finally invert the neutral rotation:
\[
  \mathcal W^3_\mu
  =
  \frac{g_2Z_\mu+g_1A_\mu}{\sqrt{g_1^2+g_2^2}},
  \qquad
  \mathcal B_\mu
  =
  \frac{-g_1Z_\mu+g_2A_\mu}{\sqrt{g_1^2+g_2^2}} .
\]
The coefficient of \(A_\mu\) in
\(g_2\mathcal W^3_\mu T^3_L+g_1\mathcal B_\mu Y\) is
\[
  \frac{g_1g_2}{\sqrt{g_1^2+g_2^2}}\,(T^3_L+Y).
\]
This proves the displayed electromagnetic coupling and identifies the
massless gauge field with the connection for the unbroken stabilizer.
\end{proof}

This is the nonabelian mass-matrix construction of
Chapter~\ref{ch:vol2-classical-yang-mills-matter} applied to the
\((\mathbf2)_{1/2}\) scalar representation.  The three coordinates
\(\chi_1,\chi_2,\chi_3\) are gauge-orbit coordinates in a local chart; in an
\(R_\xi\) gauge they remain as unphysical fields paired by BRST, while in a
unitary-gauge chart they are removed locally.

\section{Gauge-Anomaly Cancellation}

\begin{definition}[Gauge-anomaly coefficients for the local datum]
\label{def:sm-local-anomaly-coefficients}
For a collection of left-handed Weyl fermions in representations \(R_i\) of
\(\mathfrak g_{\rm SM}\), the perturbative gauge-anomaly tensor is the
symmetric cubic functional
\[
  \mathcal A(X,Y,Z)
  =
  \sum_i \operatorname{tr}_{R_i}\bigl(X\{Y,Z\}\bigr),
  \qquad
  X,Y,Z\in\mathfrak g_{\rm SM},
\]
with spectator multiplicities included in the trace.  For a mixed
\(G^2U(1)_Y\) component, this reduces to the sum of hypercharge times the
quadratic index of the \(G\)-representation.  For the mixed
gravitational-\(U(1)_Y\) anomaly, the coefficient is the sum of hypercharges
over all left-handed Weyl components, again including spectator
multiplicities.
\end{definition}

\begin{proposition}[SM anomaly cancellation]
\label{prop:sm-generation-anomaly-cancellation}
For the one-generation representation of
Definition~\ref{def:sm-one-generation-left-weyl}, all perturbative local
gauge anomalies and the mixed gravitational-hypercharge anomaly vanish.  The
finite \(SU(2)_L\) anomaly also vanishes generation by generation.
\end{proposition}

\begin{proof}
The calculation uses only the representation table and the fact that all
fermions have been written as left-handed Weyl fields.  For the pure
\(SU(3)_c^3\) anomaly, the two weak components of \(Q_L\) contribute two
fundamental color representations, while \(u^c\) and \(d^c\) contribute two
antifundamentals.  Since the cubic index changes sign under complex
conjugation,
\[
  2A(\mathbf3)-A(\mathbf3)-A(\mathbf3)=0 .
\]
The pure \(SU(2)_L^3\) perturbative anomaly vanishes because the symmetric
cubic invariant \(d^{abc}\) of \(\mathfrak{su}(2)\) is zero.  Mixed
\(SU(3)_cU(1)_Y^2\) and \(SU(2)_LU(1)_Y^2\) components vanish because they
contain \(\operatorname{tr}_{R}(T^a)\) for a semisimple generator.

For \(SU(3)_c^2U(1)_Y\), one generation gives
\[
  \mathcal A[SU(3)^2U(1)_Y]
  =
  2\left(\frac16\right)T(\mathbf3)
  +\left(-\frac23\right)T(\overline{\mathbf3})
  +\left(\frac13\right)T(\overline{\mathbf3})
  =0 .
\]
Here \(T(\overline{\mathbf3})=T(\mathbf3)\), and the factor \(2\) is the weak
doublet multiplicity.  For \(SU(2)_L^2U(1)_Y\),
\[
  \mathcal A[SU(2)^2U(1)_Y]
  =
  3\left(\frac16\right)T(\mathbf2)
  +\left(-\frac12\right)T(\mathbf2)
  =0 .
\]
Here the factor \(3\) is the color multiplicity of \(Q_L\).  For the cubic
hypercharge anomaly, include the full dimensions of the nonabelian
representations:
\[
\begin{aligned}
  \mathcal A[U(1)_Y^3]
  &=
  6\left(\frac16\right)^3
  +3\left(-\frac23\right)^3
  +3\left(\frac13\right)^3
  +2\left(-\frac12\right)^3
  +1^3\\
  &=0 .
\end{aligned}
\]
The mixed gravitational-hypercharge anomaly coefficient is
\[
  6\left(\frac16\right)
  +3\left(-\frac23\right)
  +3\left(\frac13\right)
  +2\left(-\frac12\right)
  +1
  =0 .
\]
Finally, the finite \(SU(2)\) anomaly of
Chapter~\ref{ch:volume-ii-21} is absent because each
generation contains four left-handed \(SU(2)_L\) doublets: three colored
quark doublets and one lepton doublet.  The number of doublets is therefore
even.  These equalities are finite representation-theoretic constraints on
the regulated fermion determinant line; they are not consequences of
perturbative renormalizability alone.
\end{proof}

\section{The Hybrid Definition}

\begin{definition}[Hybrid Standard Model datum]
\label{def:hybrid-standard-model-datum}
The Standard Model is used in this monograph as the following hybrid datum,
not as a single completed constructive theorem for a four-dimensional
continuum QFT:
\[
  \mathfrak D_{\rm SM}
  =
  \bigl(
    \mathfrak D_{\rm QCD}^{\rm lat},
    \mathfrak D_{\rm EW}^{\Lambda},
    \mathfrak D_{\rm weak\,EFT},
    \mathcal M
  \bigr).
\]
Here \(\mathfrak D_{\rm QCD}^{\rm lat}\) is the finite-lattice
gauge-invariant regularization of the \(SU(3)_c\) sector with quark matter
and a continuum-limit problem of the kind stated in
Hypothesis~\ref{hyp:lattice-yang-mills-continuum-limit};
\(\mathfrak D_{\rm EW}^{\Lambda}\) is the electroweak chiral gauge theory
with Higgs field treated as a finite-cutoff effective theory below a scale
\(\Lambda\); \(\mathfrak D_{\rm weak\,EFT}\) is the tower of lower-energy
effective theories obtained by integrating out massive weak fields and heavy
flavors; and \(\mathcal M\) is the matching data fixing parameters and local
operators between overlapping regimes.
\end{definition}

\begin{proposition}[Logical layers of the hybrid datum]
\label{prop:sm-hybrid-logical-layers}
The datum of Definition~\ref{def:hybrid-standard-model-datum} separates
finite-regulator definition, perturbative effective definition, and
observable-specific matching.  These layers cannot be replaced by one another
without adding an extra theorem or hypothesis.
\end{proposition}

\begin{proof}
First, the QCD sector has a precise finite nonperturbative regulator by
compact lattice link variables, Haar measure, and finite Grassmann algebra
for quarks.  The existence of the four-dimensional continuum limit with the
desired physical spectrum is still a theorem-level open problem, but the
finite-regulator object and its gauge-invariant observables are mathematically
well-defined.

Second, the electroweak sector is a perturbatively renormalizable chiral gauge
EFT.  Its perturbative consistency uses anomaly cancellation, BRST or BV
Ward identities, and the local power-counting analysis of the BEH phase.  A
standalone nonperturbative continuum construction is not supplied by these
facts.  The hypercharge \(U(1)\) coupling is not asymptotically free, and the
four-dimensional scalar quartic sector lies on the same marginal-triviality
boundary discussed in Chapters~\ref{ch:scalar-path-integrals-green-functions} and
\ref{ch:volume-ii-ising-universality}.  Moreover, a fully satisfactory
interacting lattice construction of a four-dimensional chiral gauge theory
with the Standard Model representation content is a separate problem, related
to the determinant-line and chiral-lattice issues in
Chapter~\ref{ch:constructive-lattice-fermions-chiral-symmetry}.

Third, physical predictions are extracted in regimes where a controlled
object is available.  High-energy short-distance processes use perturbative
renormalization and factorization with explicitly stated infrared-safe
observables.  Low-energy hadronic weak processes use QCD matrix elements
defined by the QCD regulator together with weak effective operators matched at
the electroweak scale.  Precision electroweak observables use the
finite-order perturbative EFT, with the cutoff dependence absorbed into
renormalized parameters and higher-dimensional operators.
\end{proof}

\begin{definition}[Standard Model prediction in the hybrid framework]
\label{def:sm-hybrid-prediction}
A Standard Model prediction for an observable \(O\) is certified in this
chapter only after the following data are specified:
\begin{enumerate}
\item the energy regime and the regulator or EFT used in that regime;
\item the local operator or detector functional representing \(O\);
\item the matching map from the adjacent higher-energy description, including
      the scheme and normalization of all operator coordinates;
\item the perturbative, lattice, or combined expansion parameter and the
      associated remainder or systematic-error class;
\item the positivity, unitarity, gauge-invariance, and infrared-safety
      statements needed for the observable.
\end{enumerate}
\end{definition}

\section{Fermi Theory as a Matched Low-Energy Layer}

\begin{proposition}[Tree-level matching to Fermi theory]
\label{prop:sm-fermi-tree-matching}
Let \(J^\pm_\mu\) be the charged left-handed weak currents normalized so that
the massive charged-vector part of the electroweak EFT is
\[
  \mathcal L_{\rm cc}
  =
  -\frac{g_2}{\sqrt2}\bigl(J^+_\mu W^{+\mu}+J^-_\mu W^{-\mu}\bigr)
  +m_W^2 W^+_\mu W^{-\mu}
  +\cdots ,
\]
where the dots include the \(W^\pm\) kinetic term and interactions.  In the
low-energy EFT for external momenta \(p^2\ll m_W^2\), eliminating \(W^\pm\) at
leading order gives
\[
  \mathcal L_{\rm Fermi}
  =
  -\frac{g_2^2}{2m_W^2}J^+_\mu J^{-\mu}
  +O(p^2/m_W^4).
\]
Therefore, with the current convention used here,
\[
  \mathcal L_{\rm Fermi}
  =
  -\frac{4G_F}{\sqrt2}J^+_\mu J^{-\mu}
  \qquad\Longleftrightarrow\qquad
  \frac{G_F}{\sqrt2}=\frac{1}{2v^2}.
\]
\end{proposition}

\begin{proof}
At tree level and to leading order in derivatives, the charged vector is an
auxiliary variable with algebraic Lagrangian
\[
  m_W^2 W^+_\mu W^{-\mu}
  -a(J^+_\mu W^{+\mu}+J^-_\mu W^{-\mu}),
  \qquad a:=\frac{g_2}{\sqrt2}.
\]
The equations obtained by varying \(W^+\) and \(W^-\) independently are
\[
  m_W^2W^-_\mu=aJ^-_\mu,\qquad
  m_W^2W^+_\mu=aJ^+_\mu .
\]
Substitution gives
\[
  \left(\frac{a^2}{m_W^2}
  -2\frac{a^2}{m_W^2}\right)J^+_\mu J^{-\mu}
  =
  -\frac{g_2^2}{2m_W^2}J^+_\mu J^{-\mu}.
\]
Including the kinetic operator replaces \(m_W^{-2}\) by the low-momentum
expansion of the massive propagator,
\[
  \frac1{m_W^2-p^2}
  =
  \frac1{m_W^2}+O(p^2/m_W^4),
\]
in the local derivative expansion.  The mass formula from
Proposition~\ref{prop:sm-electroweak-stabilizer-masses} gives
\[
  \frac{g_2^2}{2m_W^2}=\frac{2}{v^2}.
\]
Comparing with
\(-4G_FJ^+_\mu J^{-\mu}/\sqrt2\) proves
\(G_F/\sqrt2=1/(2v^2)\).
\end{proof}

The equality is a tree-level matching statement in the electroweak EFT.  QCD
renormalization of hadronic weak operators, electromagnetic corrections, CKM
factors, and matrix elements of the resulting four-fermion operators are
separate parts of the low-energy hybrid datum.

\begin{openproblem}[Constructive Standard Model]
\label{op:constructive-standard-model}
Construct a continuum local QFT realizing the Standard Model representation
content, gauge symmetry, anomaly cancellation, BEH phase, QCD confinement
data, and weak interactions with positivity, locality, covariance, and a
controlled relation to the observed parameter regime.  A solution may be a
single nonperturbative construction or a theorem explaining why the hybrid
datum above has a universal continuum domain for the specified observables.
No such theorem is assumed in this monograph.
\end{openproblem}
